\section{Evidence design and numerical implementation}
\label{sec:numerical_dynamic}

\subsection{Base history, exact generator and propagation}

The repository verification seed uses the complete 18-slip closure in Section~\ref{sec:base_model_dynamic}, a slip-gradient length $\ell_\chi=0.25~\mu\mathrm m$ and a Nye length $\ell_N=1~\mu\mathrm m$. Its grouped constraints are informed by titanium literature, but the card is not identified for a particular batch or specimen. The base history is integrated independently and then stored at 129 accepted states on $[0.5,140]~\mu\mathrm s$. At each checkpoint, the continuous constitutive Jacobian in Eq.~\eqref{eq:ft_constitutive_tangent_contract} is assembled while the local active-set and rate-sign branch is locked. A checkpoint is rejected if a derivative probe changes that branch or if
\begin{equation}
 \operatorname{cond}_2(\boldsymbol K_{\zeta\zeta})\ge10^{14},
 \qquad
 \operatorname{cond}_2(\boldsymbol R)\ge10^{12},
 \qquad\text{or}\qquad
 \operatorname{cond}_2(\boldsymbol C)\ge10^{12}.
 \label{eq:dyn_conditioning_gates}
\end{equation}
The algebraic microfield is eliminated, the generator in Eq.~\eqref{eq:dyn_generator_matrix} is assembled and the block reconstruction residual is recorded. Adjacent checkpoints are additionally scanned for yield-activity and signed-rate changes. The resulting generator path is treated as piecewise smooth. Every detected transition interval is propagated using left-continuous, right-continuous, linear-midpoint and secant-split conventions; their spread is reported as a branch-convention sensitivity rather than hidden inside the interpolation error.  For each yield event, the secant event rate and the one-sided thermally activated threshold-rate jump are also evaluated.  The released propagation uses identity saltation only after those events pass the transversality and relative-jump gates associated with Eq.~\eqref{eq:dyn_saltation_matrix}.

Between stored states $t_i$ and $t_{i+1}$, the fixed-scale generator is interpolated linearly. Four exponential-midpoint substeps are used per interval:
\begin{align}
 \widehat{\boldsymbol A}_{ij}
 &=(1-f_j)\widehat{\boldsymbol A}_i+f_j\widehat{\boldsymbol A}_{i+1},
 &f_j&=(j+1/2)/4,\\
 \boldsymbol\Phi_{ij+1}
 &=\exp(h_i\widehat{\boldsymbol A}_{ij})\boldsymbol\Phi_{ij},
 &h_i&=(t_{i+1}-t_i)/4.
 \label{eq:ft_midpoint_product}
\end{align}
The reference history therefore uses 512 ordered exponentials. Base-history refinement and propagation-substep refinement are audited separately.

The diagonal coordinate scale $\boldsymbol D_q$ is fixed before any spectrum or singular vector is inspected. Representative scales are $1~\mu\mathrm m$ for displacement, $0.1~\mathrm{m\,s^{-1}}$ for velocity, $1~\mathrm K$ for temperature, $2\times10^{-4}$ for plastic-chart and signed-slip coordinates, $1.6\times10^{12}~\mathrm{m^{-2}}$ for mobile density and $10^{10}~\mathrm{m^{-2}}$ for dipole density. Input and output selectors retain temperature and the 62 local coordinates. Displacement and velocity remain coupled inside the propagator but are not directly optimized. The declared threshold is $\mathcal G_c=e$.

Two supplementary metrics test the dependence on coordinate normalization. A recoverable-energy metric $\boldsymbol W_E(t)$ is assembled from the second variations of elastic strain energy, kinetic energy, thermal availability and condensed micromorphic curvature. Because the present stored dislocation contribution is affine in density, $\boldsymbol W_E$ is positive semidefinite rather than positive definite.  We therefore report a gain restricted to its positive input/output eigenspaces and separately test the condition $\boldsymbol\Phi\ker\boldsymbol W_{E,\mathrm{in}}\subseteq\ker\boldsymbol W_{E,\mathrm{out}}$.  Only if that condition holds may the restricted map be interpreted as an operator on quotient spaces. A second metric $\boldsymbol W_O$ treats the registered coordinate scales as independent numerical observation standard deviations. It is explicitly a numerical covariance surrogate because no synchronized DIC--thermometry--EBSD covariance is available. Both metrics are evaluated at the released direction and wavenumber; no energy-metric direction--wavenumber re-optimization is claimed.

The initial time $t_0$ is also part of the input--output definition, not a numerical nuisance parameter. Fixed-candidate gains are recomputed from $t_0=0.50$, $0.75$ and $1.00~\mu\mathrm s$ at both horizons and additionally from $2.00~\mu\mathrm s$ at the terminal horizon, thereby separating pre-peak history from late-window amplification.

\subsection{Independent descriptor audit and classical anchor}

The quadratic pencil is independently converted to an 87-coordinate first-order generalized eigenproblem with state ordering $[\widehat{\boldsymbol u},\widehat{\boldsymbol v},\widehat T,\widehat{\boldsymbol\zeta},\widehat{\boldsymbol a}]$. Generalized Schur decomposition must return 69 finite roots and 18 infinite roots. The finite roots are matched to eigenvalues of the reduced 69-state generator, and both modal backward error and microfield reconstruction residual are evaluated. No small mass or relaxation time is assigned to the algebraic coordinates.

Separately, the scalar Bai coefficient/root implementation is checked against exact-dyadic sign and zero tests and targeted regression cases. This establishes an executable classical anchor. It is not used as evidence of a strict analytical HCP-to-Bai reduction; that reduction remains a separate open gate.

\subsection{Central finite-time/frozen-spectrum discrimination}

For each fixed branch $(k,\boldsymbol n)$, define the frozen spectral abscissa of the same scaled 69-state generator as
\begin{equation}
 \alpha(t;k,\boldsymbol n)
 =\max_j\operatorname{Re}\lambda_j
 [\widetilde{\boldsymbol A}_{\rm dyn}(t;k,\boldsymbol n)]
 \label{eq:ft_spectral_abscissa}
\end{equation}
and its accumulated prediction
\begin{equation}
 \mathcal I_\alpha(t;k,\boldsymbol n)
 =\int_{t_0}^{t}\alpha(\xi;k,\boldsymbol n)\,\mathrm d\xi .
 \label{eq:ft_frozen_integral}
\end{equation}
The primary like-for-like finite-time quantity is the full-state log gain
\begin{equation}
 \mathcal L_{\rm FT}^{\rm full}(t;k,\boldsymbol n)
 =\log\|\widetilde{\boldsymbol\Phi}
 (t,t_0;k,\boldsymbol n)\|_2,
 \label{eq:ft_log_gain}
\end{equation}
and the discrimination statistic is
\begin{equation}
 \Delta_{\rm FT}^{\rm full}
 =\mathcal L_{\rm FT}^{\rm full}-\mathcal I_\alpha .
 \label{eq:ft_discrimination}
\end{equation}
Equations~\eqref{eq:ft_spectral_abscissa}--\eqref{eq:ft_discrimination}
use the same complete 69-state space, fixed scaling and identity input/output
selectors. The projected localization gain in
Eq.~\eqref{eq:dyn_observable_gain} is reported separately and is not
subtracted from the full-state $\alpha$. The selected 63-coordinate subspace
has nonzero mechanical input and output coupling, so it is not a closed
instantaneous system.

As a same-norm upper-bound audit, define the numerical abscissa
\begin{equation}
 \mu_2(\widetilde{\boldsymbol A})
 =\lambda_{\max}\!\left[
 \frac{\widetilde{\boldsymbol A}
 +\widetilde{\boldsymbol A}^{\mathsf H}}{2}\right],
 \qquad
 \log\|\widetilde{\boldsymbol\Phi}(t,t_0)\|_2
 \le\int_{t_0}^{t}\mu_2(\widetilde{\boldsymbol A})\,\mathrm dt .
 \label{eq:ft_logarithmic_norm_bound}
\end{equation}
The bound is not used as a selector; it checks that the propagated full-state
gain is consistent with the declared norm.

To verify both the implementation and the sufficiency statement, let
\begin{equation}
 \widetilde{\boldsymbol A}(t)=\boldsymbol U
 \operatorname{diag}[\lambda_1(t),\ldots,\lambda_m(t)]\boldsymbol U^{\mathsf H},
 \qquad \boldsymbol U^{\mathsf H}\boldsymbol U=\boldsymbol I,
 \label{eq:ft_commuting_normal_control}
\end{equation}
where the basis is independent of time. The family is normal and pairwise commuting, and hence
\begin{equation}
 \log\|\boldsymbol\Phi(t,t_0)\|_2
 =\max_j\int_{t_0}^{t}\operatorname{Re}\lambda_j(\xi)\,\mathrm d\xi.
 \label{eq:ft_commuting_normal_identity}
\end{equation}
This equals $\int\alpha\,\mathrm dt$ if one common mode attains the pointwise maximum almost everywhere. A stationary normal control and a time-varying commuting-normal control with a fixed leader test this equality using the same exponential-midpoint/frozen-integral quadrature. A third commuting control switches its pointwise leader and must fail the equality; it demonstrates that commutation alone is insufficient. The HCP test is then applied to the optimized near-onset $x$- and $y$-like branches and to the terminal $y$-like branch.

A second stationary diagonal control applies two invariant coordinate
projections. When the frozen rate is restricted to the same invariant
subspace, projected gain and frozen accumulation coincide; when the selected
subspace excludes the full-state leader, comparison with the full-state
$\alpha$ fails by construction. This control detects precisely the
input--output mismatch avoided in Eq.~\eqref{eq:ft_discrimination}.

The source of HCP non-equivalence is diagnosed with four quantities. For
adjacent scaled generators, the normalized commutator is
\begin{equation}
 c_i=\frac{\|\widetilde{\boldsymbol A}_{i+1}
 \widetilde{\boldsymbol A}_i-
 \widetilde{\boldsymbol A}_i
 \widetilde{\boldsymbol A}_{i+1}\|_F}
 {\|\widetilde{\boldsymbol A}_{i+1}\|_F
  \|\widetilde{\boldsymbol A}_i\|_F}.
 \label{eq:ft_commutator_diagnostic}
\end{equation}
For the instantaneous leading eigenvalue, left and right eigenvectors give
\begin{equation}
 \kappa_\lambda
 =\frac{\|\boldsymbol w\|_2\|\boldsymbol v\|_2}
 {|\boldsymbol w^{\mathsf H}\boldsymbol v|},
 \qquad
 \theta_i=\cos^{-1}\!
 \frac{|\boldsymbol v_i^{\mathsf H}\boldsymbol v_{i+1}|}
 {\|\boldsymbol v_i\|_2\|\boldsymbol v_{i+1}\|_2}.
 \label{eq:ft_modal_condition_rotation}
\end{equation}
Finally, the chronological product is compared with the reverse-ordered
product and with $\exp(\int\widetilde{\boldsymbol A}\,\mathrm dt)$, the first
Magnus surrogate. These tests separate mode ill-conditioning, eigenvector
rotation and operator ordering from the mere observation that two terminal
numbers differ.

\subsection{Cross-orientation transfer test}

Three Bunge orientations, $(0^\circ,0^\circ,0^\circ)$, $(30^\circ,90^\circ,30^\circ)$ and $(90^\circ,60^\circ,30^\circ)$, are selected independently as the minimum-, median- and maximum-storage representatives of the pre-existing deterministic orientation atlas. For each orientation, its own direction--wavenumber pair is refined; the full-state propagator norm and frozen integral are then recomputed on that same pair, base path and scaling. No orientation is fitted to increase $|\Delta_{\rm FT}^{\rm full}|$. The transfer gate requires at least three orientations, modal backward errors below $10^{-8}$ and $|\Delta_{\rm FT}^{\rm full}|\ge0.1$ at either the selection or terminal horizon for every orientation.

\subsection{Direction--wavenumber search and horizon tracking}

The primary search covers $300\le k\le3\times10^5~\mathrm{m^{-1}}$ and one closed hemisphere, using the antipodal equivalence $\boldsymbol n\sim-\boldsymbol n$. A nested scrambled Sobol sequence samples the projective sphere and $\log k$ at 64, 256 and 1024 points. Exact sample axes and the independently registered onset/final branches are retained as anchors because a narrow maximum on the equator can be missed by an unanchored finite spherical design. Empirical projective-sphere fill angle and maximum $\log k$ gap are measured at every level. Eight mutually separated starts are refined independently for both the onset and terminal objectives, with every function evaluation retained. Candidates are retained only if the last two anchor-combined maxima differ by less than 1\%, at least two starts recur in the winning basin, $k^*$ remains strictly inside the registered interval, and reevaluation on the 129-state/four-substep history lies within 1\% in gain of the independently registered reference candidate. This is an \emph{anchor-assisted basin-retention and convergence audit}. It is not called a global-search certificate because the blind Sobol design fails to discover the narrow terminal basin. At each horizon, branches within
\begin{equation}
 \mathcal G(t;k,\boldsymbol n)\ge(1-\varepsilon_G)\Gamma_D(t),
 \qquad \varepsilon_G=0.02,
 \label{eq:ft_near_optimal_gate}
\end{equation}
are retained as a near-optimal set. Temporal convergence is evaluated on 65- and 129-state histories with one, two, four and eight midpoint substeps. Norm sensitivity uses the baseline scale plus independent factor-two increases and decreases of mechanical, thermal, dislocation, plastic-kinematic and all-constitutive groups, giving 11 fixed norms. Seven selector pairs test full/full, constitutive/constitutive, full/constitutive, constitutive/full, mechanical/mechanical, mechanical/constitutive and constitutive/temperature maps. The blind Sobol maximum and the anchor-combined maximum are reported separately so that anchor dependence cannot be hidden by the audit label.

The original 11-by-7 table ranks the three registered branch candidates and is
retained as a complete fixed-candidate screen in the Supplementary Material.
To test whether that restriction hides a selector-specific basin, all seven
selectors are additionally re-optimized at onset and terminal time under the
baseline norm and four adverse norms selected before re-optimization:
mechanical${}\times0.5$, thermal${}\times0.5$,
dislocation${}\times2$ and all-constitutive${}\times2$.  The reduced-order
reconnaissance is not treated as the optimization result.  Its complete set of
winners and local starts is merged with exact axes, the three registered
branches and a new 64-point scrambled Sobol design; every member of this shared
pool is evaluated on the 129-state/four-substep propagator.  For each of the 70
contracts, the two highest mutually separated reference-pool candidates are
then refined directly on that same reference objective.  A six-point
$(\theta,\phi,\log k)$ neighbourhood must not exceed the retained result by
more than $2\times10^{-4}$ in log gain; failed contracts are restarted from
the improving probe.  Search-to-reference reevaluation is therefore not
reported as reference re-optimization.  Lower-$k$ and upper-$k$ boundary
solutions are extended by two decades in the corresponding direction and are
classified as long-wave thermal or reversible high-frequency branches rather
than finite localization wavelengths.  Interior finite-band selection is
required only for maps that observe a constitutive-state localization
response; the scalar temperature output is allowed to select a long-wave
boundary branch.

For a strict same-input/output freezing baseline, divide every registered
nonuniform interval into $r$ equal subintervals and define
\begin{equation}
 \boldsymbol\Phi_{\rm PF}^{(r)}
 =\prod_{m=N_r}^{1}
 \exp\!\left[\Delta t_m\,
 \widetilde{\boldsymbol A}(t_{m-1})\right],
 \qquad
 \mathcal G_{\rm PF}^{(r)}
 =\left\|\boldsymbol P_{\rm out}
 \boldsymbol\Phi_{\rm PF}^{(r)}
 \boldsymbol P_{\rm in}^{\mathsf T}\right\|_2 .
 \label{eq:ft_piecewise_frozen_same_selector}
\end{equation}
The metric and both selectors are identical to those of the time-ordered
reference; only the generator is frozen at the left endpoint of each
subinterval.  Results are reported for $r=1,2,4,8,16$ subdivisions per
registered interval.  This is a
controlled discretization of chronology, not the incompatible scalar
$\int\alpha\,\mathrm dt$ comparison.

\subsection{Loading transfer, local parameter sensitivity and external titanium check}

Three additional paths are integrated without refitting: simple shear at
$5.0\times10^3$ and $1.3\times10^4~\mathrm{s^{-1}}$ from 293.15 K, and the
baseline $1.0\times10^4~\mathrm{s^{-1}}$ path from 373.15 K. Direction and wavenumber are
re-optimized for the constitutive/constitutive map at onset and terminal time
using shared Sobol reconnaissance plus continuation from the baseline basin.
These rate/temperature cases test transfer, not specimen prediction.

Local sensitivity uses symmetric multiplicative perturbations of $\pm20\%$
for the initial mobile density $\rho_{m0}$, mean-free-path coefficient
$K_\Lambda$, Nye length $\ell_N$ and full-rank slip-gradient length
$\ell_\chi$.  The same branch is continued from the baseline and
direction--wavenumber optimization is repeated.  Central logarithmic
sensitivities of gain and wavenumber are
\begin{equation}
 S_p^Y=\frac{\log Y(1.2p)-\log Y(0.8p)}
 {\log(1.2)-\log(0.8)},
 \qquad Y\in\{\mathcal G,k^*\}.
 \label{eq:ft_local_log_sensitivity}
\end{equation}
This is a branch-local derivative audit, not global uncertainty
quantification.

Finally, the $1.3\times10^4~\mathrm{s^{-1}}$ calculation is compared with the
synchronized shear-stress and infrared-temperature landmarks measured for
commercial-purity Grade-II titanium by Guo et al.\ \cite{GuoEtAl2019PRL}.
Peak and shear-band times are read from their Fig.~1; the temperature-rise
interval at peak stress and localized-temperature interval are taken from the
paper text.  Because specimen geometry, texture, damage and localization are
not matched, this is an external timing/scale check with explicit failures,
not parameter identification.

\subsection{Fixed-base coupling interventions}

To distinguish coordinate participation from dependence on constitutive coupling, the 69-state generator is partitioned into mechanical/inertial, thermal, dislocation and plastic-kinematic groups. Within-group blocks and the base history are held fixed. Six cross-group edge families are switched on or off, and the leading singular value and optimal vectors are recomputed for all $2^6=64$ coalitions. The value function is the maximum finite-time observable log gain. Exact Shapley allocation is used to summarize edge dependence, while grouped removals (no inertial, no thermal or no dislocation cross-links) provide transparent counterfactual checks. These operations identify dependence of the \emph{linearized fixed-base model} on selected blocks; they are not experimental causal effects, dissipation fractions or unique structural mechanisms.

Gradient controls are reported in the same figure but under a separate contract. They use an assembled frozen quasi-static HCP operator, extend the sampled range to $k=10^{12}~\mathrm{m^{-1}}$, and test removal of slip-gradient and Nye contributions. They do not supply nonlinear widths or a calibrated high-$k$ theorem, so no quantitative comparison between the dynamic coupling-intervention panel and the quasi-static gradient panel is made.

\subsection{Singular-vector nonlinear transport and material-width gates}

The first nonlinear gate uses the actual leading right singular vectors of the
re-optimized baseline constitutive/constitutive map at $1.4375~\mu\mathrm s$
and $140~\mu\mathrm s$.  Candidate direction and wavenumber and the coordinate
metric are held at the direct-reference V4 winners obtained on the
129-state/four-substep objective.  They are not inherited from the lower-resolution
reconnaissance.  The fixed-$\beta$, passive-ledger base path
is independently refined from 129 to 513 states, and the singular vectors are
recomputed on that refined map before nonlinear transport.  Dense-history
convergence compares nested 257- and 513-state maps with four midpoint
substeps, followed by a one-to-four-substep check on the 513-state map.  If
the 257-to-513 relative log-gain change exceeds 2\%, a 1025-state map is
mandatory and convergence is judged from the 513-to-1025 change under the
same four-substep contract.  The released 129-state gain is retained as a
coarse-history screen; a finer value supersedes it only after this density
gate closes, without changing the registered direction, wavenumber or metric.
Each V4 vector is imposed as a
real periodic Fourier mode
with 16 cells per wavelength and two factor-two amplitudes.  The first replay
retains nonnegative modes 0--3 and is therefore labelled a Galerkin transport
check, not a complete Fourier calculation.  The terminal
amplitude is additionally limited so that its linearly predicted final
dimensionless norm does not exceed $10^{-3}$. The nonlinear periodic
momentum--heat--micromorphic--crystal-state residual is advanced with the
registered tangent used only as an integrating factor. Acceptance requires
less than 2\% error in gain, less than 3\% error in the complete complex
Fourier output vector, less than 2\% factor-two gain variation, less than 1\%
nonretained Fourier energy within the declared subspace and a power-partition
ledger residual below $10^{-10}$.  Passing these Galerkin checks does not
override the independent full-grid gate.

That gate separates discretization convergence from nonlinear spectral
closure.  On the 16-cell grid, a raw mode-0--8 calculation is retained as a
negative control, while an otherwise identical mode-0--7 calculation removes
only the even-grid Nyquist coordinate.  A dealiased ladder then retains
nonnegative modes $0,\ldots,\lfloor N/3\rfloor$ on $N=16$, 32 and 64 cells.
Spatial convergence requires the 32-to-64-cell gain change to be below 1\% at
both horizons.  Time convergence independently compares 257-state histories
with one, two and four substeps against a 513-state one-substep reference.
Integration-format sensitivity compares exponential midpoint with a
Lawson--Euler control and must remain below 2\%.  The primary dealiased replay
must also reproduce the linear gain within 2\% and the complete retained
complex output vector within 3\%.  Every non-finite run records the first
failed time interval, residual stage, state coordinate and signed
wavenumber.  Finally, the instantaneous nonlinear right-hand side is projected
onto discarded modes; more than 1\% discarded energy leaves nonlinear
harmonic closure open even when the resolved state and gain are converged.

A separate $0$--$2~\mu\mathrm s$ material-width bridge imposes prescribed sinusoidal or multimode disturbances with locked phases and spatial grids. It is accepted in stages. First, linear gain and phase transport must be reproduced. Second, multiple input amplitudes must remain in the amplitude-linear regime for a single mode. Third, the peak and full width at half maximum (FWHM) must converge on nested grids. Finally, a candidate material width must be insensitive to seed phase/spectrum, remain resolved by at least 12 cells and emerge independently of the imposed wavelength. A failure retains the successful lower-level verification but disables the material-scale interpretation.
